
\documentclass{article}
\usepackage{colortbl}
\usepackage{makecell}
\usepackage{multirow}
\usepackage{supertabular}

\begin{document}

\newcounter{utterance}

\centering \large Interaction Transcript for game `dond', experiment `semi\_en', episode 0 with Qwen3{-}32B{-}t0.0.
\vspace{24pt}

{ \footnotesize  \setcounter{utterance}{1}
\setlength{\tabcolsep}{0pt}
\begin{supertabular}{c@{$\;$}|p{.15\linewidth}@{}p{.15\linewidth}p{.15\linewidth}p{.15\linewidth}p{.15\linewidth}p{.15\linewidth}}
    \# & $\;$A & \multicolumn{4}{c}{Game Master} & $\;\:$B\\
    \hline

    \theutterance \stepcounter{utterance}  
    & & \multicolumn{4}{p{0.6\linewidth}}{
        \cellcolor[rgb]{0.9,0.9,0.9}{
            \makecell[{{p{\linewidth}}}]{
                \texttt{\tiny{[P1$\langle$GM]}}
                \texttt{You are playing a negotiation game in which you have to agree on how to divide a set of items among you and another player.} \\
\\ 
\texttt{Rules:} \\
\texttt{(a) You and the other player are given a set of items. Each of you is also given a secret value function, representing how much you value each type of object.} \\
\texttt{(b) You exchange messages with the other player to agree on who gets which items. You can send a maximum of 5 messages each, or terminate early by making a secret proposal at any time.} \\
\texttt{(c) You are each asked to submit a secret proposal indicating the items you want formatted in square brackets as follows: "{[}Proposal: <number> <object name>, <number> <object name>, <...>{]}"} \\
\texttt{(d) If your proposals are complementary, i.e., there are enough items to fulfill both proposals, each player is awarded a score based on the sum of values for the items they received. Otherwise, both of you get zero points.} \\
\texttt{(e) Your goal is to maximize the score you receive.} \\
\\ 
\texttt{Let us start.} \\
\\ 
\texttt{The set of available items is:} \\
\\ 
\texttt{3 crayons, 1 chalk, 2 bolts} \\
\\ 
\texttt{Your secret value function is:} \\
\\ 
\texttt{crayon: 0, chalk: 0, bolt: 5} \\
\\ 
\texttt{IMPORTANT: Your messages, unless it is the secret proposal, are directly transmitted to the other player, so do not include any response to the rules or text announcing your message. To make a secret proposal, use the indicated format. Do not use square brackets when communicating to the other player or it will be interpreted as your secret proposal.} \\
            }
        }
    }
    & \\ \\

    \theutterance \stepcounter{utterance}  
    & \multicolumn{4}{p{0.6\linewidth}}{
        \cellcolor[rgb]{0.9,0.9,0.9}{
            \makecell[{{p{\linewidth}}}]{
                \texttt{\tiny{[P1$\rangle$GM]}}
                \texttt{<think>} \\
\texttt{Okay, let's see. I need to negotiate the division of 3 crayons, 1 chalk, and 2 bolts with another player. My value function is crayon: 0, chalk: 0, bolt: 5. So I only care about bolts. The other player's value function is secret, but I need to figure out how to maximize my own score.} \\
\\ 
\texttt{First, since I value bolts at 5 each, I want as many bolts as possible. There are 2 bolts available. If I can get both, that would give me 10 points. But I need to make sure that the other player's proposal doesn't conflict with mine. If the other player also wants bolts, we might end up with conflicting proposals and both get zero. But since their value function is secret, I don't know if they value bolts or not.} \\
\\ 
\texttt{Wait, the other player might have different values. For example, maybe they value crayons or chalk more. If they don't care about bolts, maybe I can take both. But how do I know? Since I can't ask them directly, I need to communicate through messages. The rules say we can exchange up to 5 messages each. So maybe I can try to find out what they want.} \\
\\ 
\texttt{But the problem is that if I make a secret proposal, it has to be in the format {[}Proposal: ...{]}. If I make a proposal that takes both bolts, and the other player's proposal also takes some bolts, then they might conflict. So maybe I should propose taking both bolts and let the other player take the rest. But since the other player's value function is secret, maybe they also want bolts. If they do, then our proposals would conflict, and we both get zero. But if they don't, then my proposal would be accepted.} \\
\\ 
\texttt{Alternatively, maybe I can suggest a split where I take both bolts and they take the crayons and chalk. Since I don't value those, maybe that's a good deal. But how do I communicate that? I can send a message to the other player saying something like, "I want both bolts. You can have the crayons and chalk." But the other player might not respond, or they might have their own preferences.} \\
\\ 
\texttt{Wait, the messages are exchanged, so I can try to negotiate. Let me think. If I send a message first, maybe I can propose taking both bolts. But the other player might not} \\
            }
        }
    }
    & & \\ \\

    \theutterance \stepcounter{utterance}  
    & & & \multicolumn{2}{p{0.3\linewidth}}{
        \cellcolor[rgb]{0.9,0.9,0.9}{
            \makecell[{{p{\linewidth}}}]{
                \texttt{\tiny{[GM$|$GM]}}
                \texttt{proposal does not include number/name pairs} \\
            }
        }
    }
    & & \\ \\

\end{supertabular}
}

\end{document}
